% GENERATED FILE. Edit canonical lessons, then run npm run generate:books.
% canonical-source-hash: fnv1a64:5e0dc2947e9d5dc4
% canonical-lessons: BN-C347-doptor, BN-C347-pourosobha, BN-C347-ponchayet, BN-C347-niyog-kora, BN-C347-jari-kora

\chapter{Public Offices}
\label{ch:bn-public-offices}
\hlchaptermodality{347}

\begin{chapteropening}
\textbf{By the end of this chapter:} \emph{I can say which office deals with what.}

Complete the last lesson of chapter 347: I can say which office deals with what.
\end{chapteropening}

\section[\texorpdfstring{\bn{দপ্তর} (dôptôr)}{দপ্তর (dôptôr)}]{\bn{দপ্তর} (dôptôr) — a government department; an office}

\label{lesson:BN-C347-doptor}



\begin{quote}
Before the new one: say the Bengali for to resign, then the Bengali for to take an oath.
\end{quote}

\subsection*{You'll want to know: \bn{দপ্তর}}
\textbf{\bn{দপ্তর}} — \emph{dôptôr} — \textquotedblleft{}a government department; an office\textquotedblright{}.

\begin{quote}
\bn{আমি} \bn{আবেদনটা} \bn{দপ্তরে} \bn{জমা} \bn{দিলাম।}
\end{quote}

\begin{quote}
\emph{āmi ābedônṭā dôptôre jômā dilām.}
\end{quote}

\begin{quote}
I handed in the application at the department.
\end{quote}

\begin{cousinweb}
\textbf{\bn{দপ্তর}} (\emph{dôptôr}) is Persian \emph{daftar}, a register and then the office that keeps it, usually traced back through Arabic to Greek \emph{diphthérā}, a prepared hide. Bengali sounds it with p.
\end{cousinweb}

\begin{grammarlens}[title={What you've built}]
One more word for giving reasons and opinions.
\end{grammarlens}

\subsection*{Your turn}
\begin{itemize}
\raggedright
  \item \emph{Say it:} \emph{dôptôr}
  \item \emph{Say it:} \emph{dôptôr}, once more
  \item \emph{Say it:} \emph{āmi ābedônṭā dôptôre jômā dilām.} — I handed in the application at the department.
  \item \emph{Recall:} say the Bengali for to resign, then the Bengali for to take an oath, then say \emph{dôptôr} again
\end{itemize}

\begin{tcolorbox}[breakable,colback=teal!4,colframe=teal!35!black,title={Before you move on}]
What does \bn{দপ্তর} mean? (\textquotedblleft{}A government department; an office\textquotedblright{}.) Say it once more.
\end{tcolorbox}
\section[\texorpdfstring{\bn{পৌরসভা} (pourôsôbhā)}{পৌরসভা (pourôsôbhā)}]{\bn{পৌরসভা} (pourôsôbhā) — a municipality; a town council}

\label{lesson:BN-C347-pourosobha}



\begin{quote}
Before the new one: say the Bengali for to take an oath, then the Bengali for a government department.
\end{quote}

\subsection*{You'll want to know: \bn{পৌরসভা}}
\textbf{\bn{পৌরসভা}} — \emph{pourôsôbhā} — \textquotedblleft{}a municipality; a town council\textquotedblright{}.

\begin{quote}
\bn{পৌরসভা} \bn{পাড়ার} \bn{রাস্তায়} \bn{নতুন} \bn{বাতি} \bn{লাগাল।}
\end{quote}

\begin{quote}
\emph{pourôsôbhā pāṛār rāstāy nôtun bāti lāgālô.}
\end{quote}

\begin{quote}
The town council put up new lamps in the neighbourhood's streets.
\end{quote}

\begin{cousinweb}
\textbf{\bn{পৌরসভা}} (\emph{pourôsôbhā}) is Sanskrit \emph{paura}, a townsman (Monier-Williams), with \bn{সভা}, an assembly: the town's council.
\end{cousinweb}

\begin{grammarlens}[title={What you've built}]
One more word for giving reasons and opinions.
\end{grammarlens}

\subsection*{Your turn}
\begin{itemize}
\raggedright
  \item \emph{Say it:} \emph{pourôsôbhā}
  \item \emph{Say it:} \emph{pourôsôbhā}, once more
  \item \emph{Say it:} \emph{pourôsôbhā pāṛār rāstāy nôtun bāti lāgālô.} — The town council put up new lamps in the neighbourhood's streets.
  \item \emph{Recall:} say the Bengali for to take an oath, then the Bengali for a government department, then say \emph{pourôsôbhā} again
\end{itemize}

\begin{tcolorbox}[breakable,colback=teal!4,colframe=teal!35!black,title={Before you move on}]
What does \bn{পৌরসভা} mean? (\textquotedblleft{}A municipality; a town council\textquotedblright{}.) Say it once more.
\end{tcolorbox}
\section[\texorpdfstring{\bn{পঞ্চায়েত} (pônchāyet)}{পঞ্চায়েত (pônchāyet)}]{\bn{পঞ্চায়েত} (pônchāyet) — a village council}

\label{lesson:BN-C347-ponchayet}



\begin{quote}
Before the new one: say the Bengali for a government department, then the Bengali for a town council.
\end{quote}

\subsection*{You'll want to know: \bn{পঞ্চায়েত}}
\textbf{\bn{পঞ্চায়েত}} — \emph{pônchāyet} — \textquotedblleft{}a village council\textquotedblright{}.

\begin{quote}
\bn{পঞ্চায়েতের} \bn{সভায়} \bn{গ্রামের} \bn{সবাই} \bn{এল।}
\end{quote}

\begin{quote}
\emph{pônchāyeter sôbhāy grāmer sôbāi elô.}
\end{quote}

\begin{quote}
Everyone in the village came to the council meeting.
\end{quote}

\begin{cousinweb}
\textbf{\bn{পঞ্চায়েত}} (\emph{pônchāyet}) has a first part usually taken as Sanskrit \emph{pañca}, five (CDIAL 7655 lists Bengali pā̃c), after an old village council of five elders; the rest of the word is not settled, and the whole word is in neither CDIAL nor Monier-Williams. Today it is the elected council of a village or a group of villages.
\end{cousinweb}

\begin{grammarlens}[title={What you've built}]
One more word for giving reasons and opinions.
\end{grammarlens}

\subsection*{Your turn}
\begin{itemize}
\raggedright
  \item \emph{Say it:} \emph{pônchāyet}
  \item \emph{Say it:} \emph{pônchāyet}, once more
  \item \emph{Say it:} \emph{pônchāyeter sôbhāy grāmer sôbāi elô.} — Everyone in the village came to the council meeting.
  \item \emph{Recall:} say the Bengali for a government department, then the Bengali for a town council, then say \emph{pônchāyet} again
\end{itemize}

\begin{tcolorbox}[breakable,colback=teal!4,colframe=teal!35!black,title={Before you move on}]
What does \bn{পঞ্চায়েত} mean? (\textquotedblleft{}A village council\textquotedblright{}.) Say it once more.
\end{tcolorbox}
\section[\texorpdfstring{\bn{নিয়োগ} \bn{করা} (niyog kôrā)}{নিয়োগ করা (niyog kôrā)}]{\bn{নিয়োগ} \bn{করা} (niyog kôrā) — to appoint; to take on (staff)}

\label{lesson:BN-C347-niyog-kora}



\begin{quote}
Before the new one: say the Bengali for a town council, then the Bengali for a village council.
\end{quote}

\subsection*{You'll want to know: \bn{নিয়োগ} \bn{করা}}
\textbf{\bn{নিয়োগ} \bn{করা}} — \emph{niyog kôrā} — \textquotedblleft{}to appoint; to take on (staff)\textquotedblright{}.

\begin{quote}
\bn{কোম্পানি} \bn{দশজন} \bn{নতুন} \bn{কর্মী} \bn{নিয়োগ} \bn{করল।}
\end{quote}

\begin{quote}
\emph{kompāni dôshjôn nôtun kôrmī niyog kôrlô.}
\end{quote}

\begin{quote}
The company took on ten new workers.
\end{quote}

Past (he/she …): \bn{তিনি} \bn{একজন} \bn{নতুন} \bn{সচিব} \bn{নিয়োগ} \bn{করলেন}, \emph{tini ekjôn nôtun sôchib niyog kôrlen}.

\begin{cousinweb}
\textbf{\bn{নিয়োগ} \bn{করা}} (\emph{niyog kôrā}) is Sanskrit \emph{niyoga}, a charge or appointed task (Monier-Williams), with \bn{করা}: to put someone in a post.
\end{cousinweb}

\begin{grammarlens}[title={What you've built}]
One more word for giving reasons and opinions.
\end{grammarlens}

\subsection*{Your turn}
\begin{itemize}
\raggedright
  \item \emph{Say it:} \emph{niyog kôrā}
  \item \emph{Say it:} \emph{niyog kôrā}, once more
  \item \emph{Say it:} \emph{kompāni dôshjôn nôtun kôrmī niyog kôrlô.} — The company took on ten new workers.
  \item \emph{Recall:} say the Bengali for a town council, then the Bengali for a village council, then say \emph{niyog kôrā} again
\end{itemize}

\begin{tcolorbox}[breakable,colback=teal!4,colframe=teal!35!black,title={Before you move on}]
What does \bn{নিয়োগ} \bn{করা} mean? (\textquotedblleft{}To appoint; to take on (staff)\textquotedblright{}.) Say it once more.
\end{tcolorbox}
\section[\texorpdfstring{\bn{জারি} \bn{করা} (jāri kôrā)}{জারি করা (jāri kôrā)}]{\bn{জারি} \bn{করা} (jāri kôrā) — to issue an order; to issue (a notice)}

\label{lesson:BN-C347-jari-kora}



\begin{quote}
Before the new one: say the Bengali for a village council, then the Bengali for to appoint.
\end{quote}

\subsection*{You'll want to know: \bn{জারি} \bn{করা}}
\textbf{\bn{জারি} \bn{করা}} — \emph{jāri kôrā} — \textquotedblleft{}to issue an order; to issue (a notice)\textquotedblright{}.

\begin{quote}
\bn{সরকার} \bn{একটা} \bn{নতুন} \bn{নোটিশ} \bn{জারি} \bn{করল।}
\end{quote}

\begin{quote}
\emph{sôrkār ekṭā nôtun noṭish jāri kôrlô.}
\end{quote}

\begin{quote}
The government issued a new notice.
\end{quote}

Past (he/she …): \bn{তিনি} \bn{নোটিশটা} \bn{জারি} \bn{করলেন}, \emph{tini noṭishṭā jāri kôrlen}.

\begin{cousinweb}
\textbf{\bn{জারি} \bn{করা}} (\emph{jāri kôrā}) is Persian \emph{jārī}, flowing or current, usually traced to Arabic, with \bn{করা}: to set an order flowing, to bring it into force.
\end{cousinweb}

\begin{grammarlens}[title={What you've built}]
One more word for giving reasons and opinions.
\end{grammarlens}

\subsection*{Your turn}
\begin{itemize}
\raggedright
  \item \emph{Say it:} \emph{jāri kôrā}
  \item \emph{Say it:} \emph{jāri kôrā}, once more
  \item \emph{Say it:} \emph{sôrkār ekṭā nôtun noṭish jāri kôrlô.} — The government issued a new notice.
  \item \emph{Recall:} say the Bengali for a village council, then the Bengali for to appoint, then say \emph{jāri kôrā} again
\end{itemize}

\begin{tcolorbox}[breakable,colback=teal!4,colframe=teal!35!black,title={Before you move on}]
What does \bn{জারি} \bn{করা} mean? (\textquotedblleft{}To issue an order; to issue (a notice)\textquotedblright{}.) Say it once more.
\end{tcolorbox}
